% Source: Weinberg GR, physical PDF pp. 235--243
%         (printed pp. 212--220).
% Coverage: Section 9.1, The Post-Newtonian Approximation; equations
%           (9.1.1)--(9.1.75).
% Figures/tables/footnotes: no figures or tables.
% Status: source-reviewed and compile-clean.
% Uncertainties: Weinberg's stacked perturbative-order numerals are rendered
%                as parenthesized superscripts; curvature-sign conversion is
%                documented immediately below and has been checked against
%                the target Newtonian limit.
%
% MODERNIZATION CHECK: Weinberg prints perturbative order as a numeral over
% the tensor symbol.  We use a parenthesized superscript, e.g.
% g_{\mu\nu}^{(N)}.  This is an order label, not a tensor index or power.
% MODERNIZATION CHECK: Equations (9.1.26)--(9.1.41) use the target curvature
% convention.  Thus every displayed Ricci component is the negative of the
% source expression at fixed slots, while the trace-reversed field equation
% has +8 pi G rather than the source's -8 pi G.  The metric solutions
% (9.1.53)--(9.1.75) are consequently unchanged.

\subsection{The Post-Newtonian Approximation}\label{sec:9.1}

Consider a system of particles that, like the sun and planets, are bound
together by their mutual gravitational attraction. Let
\(\bar M,\bar r,\bar v\) be typical values of the masses, separations, and
velocities of these particles. It is a familiar result of Newtonian mechanics
that the typical kinetic energy \(\bar M\bar v^2/2\) will be roughly of the
same order of magnitude as the typical potential energy
\(G\bar M^2/\bar r\), so
\begin{equation}
  \bar v^2\sim\frac{G\bar M}{\bar r}.
  \tag{9.1.1}\label{eq:9.1.1}
\end{equation}
(For instance, a test particle in a circular orbit of radius \(r\) about a
central mass \(M\) will have velocity \(v\) given in Newtonian mechanics by
the exact formula \(v^2=GM/r\).) The post-Newtonian approximation may be
described as a method for obtaining the motions of the system to one higher
power of the small parameters \(G\bar M/\bar r\) and \(\bar v^2\) than given
by Newtonian mechanics. It is sometimes referred to as an expansion in
inverse powers of the speed of light, but since in our units this speed is
unity we prefer to say that our expansion parameter is \(\bar v^2\), or
equivalently \(G\bar M/\bar r\).

We must begin by asking what we need. The equations of motion of the
particles are
\[
  \frac{\dd^2x^\mu}{\dd\tau^2}
  +\Gamma^\mu{}_{\nu\lambda}
   \frac{\dd x^\nu}{\dd\tau}\frac{\dd x^\lambda}{\dd\tau}=0.
\]
Using \(t=x^0\), we may compute the coordinate accelerations as
\[
  \begin{aligned}
  \frac{\dd^2x^i}{\dd t^2}
  &=
  \left(\frac{\dd t}{\dd\tau}\right)^{-1}
  \frac{\dd}{\dd\tau}
  \left[
    \left(\frac{\dd t}{\dd\tau}\right)^{-1}
    \frac{\dd x^i}{\dd\tau}
  \right]\\
  &=
  \left(\frac{\dd t}{\dd\tau}\right)^{-2}
  \frac{\dd^2x^i}{\dd\tau^2}
  -
  \left(\frac{\dd t}{\dd\tau}\right)^{-3}
  \frac{\dd^2t}{\dd\tau^2}\frac{\dd x^i}{\dd\tau}\\
  &=
  -\Gamma^i{}_{\nu\lambda}
   \frac{\dd x^\nu}{\dd t}\frac{\dd x^\lambda}{\dd t}
  +
  \Gamma^0{}_{\nu\lambda}
   \frac{\dd x^\nu}{\dd t}\frac{\dd x^\lambda}{\dd t}
   \frac{\dd x^i}{\dd t}.
  \end{aligned}
\]
In detail, this is
\begin{equation}
  \begin{split}
  \frac{\dd^2x^i}{\dd t^2}
  ={}&
  -\Gamma^i{}_{00}
  -2\Gamma^i{}_{0j}\frac{\dd x^j}{\dd t}
  -\Gamma^i{}_{jk}
   \frac{\dd x^j}{\dd t}\frac{\dd x^k}{\dd t}\\
  &+
  \left(
    \Gamma^0{}_{00}
    +2\Gamma^0{}_{0j}\frac{\dd x^j}{\dd t}
    +\Gamma^0{}_{jk}
     \frac{\dd x^j}{\dd t}\frac{\dd x^k}{\dd t}
  \right)\frac{\dd x^i}{\dd t}.
  \end{split}
  \tag{9.1.2}\label{eq:9.1.2}
\end{equation}
In the Newtonian approximation discussed in Section~\ref{sec:3.4}, we
treated all velocities as vanishingly small and kept only terms of first
order in the difference between \(g_{\mu\nu}\) and the Minkowski tensor
\(\eta_{\mu\nu}\), and found
\[
  \frac{\dd^2x^i}{\dd t^2}
  \simeq-\Gamma^i{}_{00}
  \simeq\frac12\partial_i g_{00}.
\]
But \(g_{00}+1\) is of order \(G\bar M/\bar r\), so the Newtonian
approximation gives \(\dd^2x^i/\dd t^2\) to order
\(G\bar M/\bar r^2\), that is, to order \(\bar v^2/\bar r\).
\emph{Therefore our objective in using the post-Newtonian approximation will
be to compute \(\dd^2x^i/\dd t^2\) to order \(\bar v^4/\bar r\).}
Inspection of Eq.~\eqref{eq:9.1.2} shows that we shall need the various
components of the affine connection to the following orders:
\begin{equation}
  \begin{aligned}
  \Gamma^i{}_{00}&:\quad \bar v^4/\bar r,&
  \Gamma^i{}_{0j}&:\quad \bar v^3/\bar r,&
  \Gamma^i{}_{jk}&:\quad \bar v^2/\bar r,\\
  \Gamma^0{}_{00}&:\quad \bar v^3/\bar r,&
  \Gamma^0{}_{0j}&:\quad \bar v^2/\bar r,&
  \Gamma^0{}_{jk}&:\quad \bar v/\bar r.
  \end{aligned}
  \tag{9.1.3}\label{eq:9.1.3}
\end{equation}

From our experience with the Schwarzschild solution, we expect that it
should be possible to find a coordinate system in which the metric tensor is
nearly equal to the Minkowski tensor, the corrections being expandable in
powers of \(G\bar M/\bar r\sim\bar v^2\). In particular, we expect
\begin{equation}
  g_{00}=-1+g_{00}^{(2)}+g_{00}^{(4)}+\cdots,
  \tag{9.1.4}\label{eq:9.1.4}
\end{equation}
\begin{equation}
  g_{ij}=\delta_{ij}+g_{ij}^{(2)}+g_{ij}^{(4)}+\cdots,
  \tag{9.1.5}\label{eq:9.1.5}
\end{equation}
\begin{equation}
  g_{i0}=g_{i0}^{(3)}+g_{i0}^{(5)}+\cdots,
  \tag{9.1.6}\label{eq:9.1.6}
\end{equation}
where \(g_{\mu\nu}^{(N)}\) denotes the term in \(g_{\mu\nu}\) of order
\(\bar v^N\). Odd powers occur in \(g_{i0}\), because \(g_{i0}\) must change
sign under the time-reversal transformation \(t\to-t\). The real
justification for these expansions will come below, when we show that they
lead to a consistent solution of the Einstein field equations.

The inverse metric is defined by
\begin{equation}
  g^{i\mu}g_{0\mu}=g^{i0}g_{00}+g^{ij}g_{j0}=0,
  \tag{9.1.7}\label{eq:9.1.7}
\end{equation}
\begin{equation}
  g^{0\mu}g_{0\mu}=g^{00}g_{00}+g^{0\ii}g_{0\ii}=1,
  \tag{9.1.8}\label{eq:9.1.8}
\end{equation}
\begin{equation}
  g^{i\mu}g_{j\mu}=g^{i0}g_{j0}+g^{ik}g_{jk}=\delta^i{}_j.
  \tag{9.1.9}\label{eq:9.1.9}
\end{equation}
We expect that
\begin{equation}
  g^{00}=-1+g^{00(2)}+g^{00(4)}+\cdots,
  \tag{9.1.10}\label{eq:9.1.10}
\end{equation}
\begin{equation}
  g^{ij}=\delta^{ij}+g^{ij(2)}+g^{ij(4)}+\cdots,
  \tag{9.1.11}\label{eq:9.1.11}
\end{equation}
\begin{equation}
  g^{i0}=g^{i0(3)}+g^{i0(5)}+\cdots.
  \tag{9.1.12}\label{eq:9.1.12}
\end{equation}
Inserting these expansions into Eqs.~\eqref{eq:9.1.7}--\eqref{eq:9.1.9}
gives, to the first nontrivial orders,
\begin{equation}
  g^{00(2)}=-g_{00}^{(2)},\qquad
  g^{ij(2)}=-g_{ij}^{(2)},\qquad
  g^{i0(3)}=g_{i0}^{(3)},\quad\text{etc.}
  \tag{9.1.13}\label{eq:9.1.13}
\end{equation}

The affine connection is obtained from
\[
  \Gamma^\mu{}_{\nu\lambda}
  =\frac12g^{\mu\rho}
  \left(
    \partial_\lambda g_{\rho\nu}
    +\partial_\nu g_{\rho\lambda}
    -\partial_\rho g_{\nu\lambda}
  \right).
\]
In computing it we must take account of the fact that the scales of distance
and time in our systems are set by \(\bar r\) and \(\bar r/\bar v\),
respectively, so
\[
  \partial_i\sim\frac1{\bar r},
  \qquad
  \partial_t\sim\frac{\bar v}{\bar r}.
\]
Using Eqs.~\eqref{eq:9.1.4}--\eqref{eq:9.1.6} and
\eqref{eq:9.1.10}--\eqref{eq:9.1.13}, the components
\(\Gamma^i{}_{00},\Gamma^i{}_{jk},\Gamma^0{}_{0\ii}\) have expansions
\begin{equation}
  \Gamma^\mu{}_{\nu\lambda}
  =\Gamma^{\mu(2)}{}_{\nu\lambda}
  +\Gamma^{\mu(4)}{}_{\nu\lambda}+\cdots,
  \quad
  (\Gamma^i{}_{00},\Gamma^i{}_{jk},\Gamma^0{}_{0\ii}),
  \tag{9.1.14}\label{eq:9.1.14}
\end{equation}
while \(\Gamma^i{}_{0j},\Gamma^0{}_{00},\Gamma^0{}_{ij}\) have expansions
\begin{equation}
  \Gamma^\mu{}_{\nu\lambda}
  =\Gamma^{\mu(3)}{}_{\nu\lambda}
  +\Gamma^{\mu(5)}{}_{\nu\lambda}+\cdots,
  \quad
  (\Gamma^i{}_{0j},\Gamma^0{}_{00},\Gamma^0{}_{ij}).
  \tag{9.1.15}\label{eq:9.1.15}
\end{equation}
Here \(\Gamma^{\mu(N)}{}_{\nu\lambda}\) denotes the term of order
\(\bar v^N/\bar r\). The components required in
Eq.~\eqref{eq:9.1.3} are explicitly
\begin{equation}
  \Gamma^{i(2)}{}_{00}=-\frac12\partial_i g_{00}^{(2)},
  \tag{9.1.16}\label{eq:9.1.16}
\end{equation}
\begin{equation}
  \Gamma^{i(4)}{}_{00}
  =-\frac12\partial_i g_{00}^{(4)}
   +\partial_t g_{i0}^{(3)}
   +\frac12g_{ij}^{(2)}\partial_j g_{00}^{(2)},
  \tag{9.1.17}\label{eq:9.1.17}
\end{equation}
\begin{equation}
  \Gamma^{i(3)}{}_{0j}
  =\frac12\left(
    \partial_jg_{i0}^{(3)}
    +\partial_tg_{ij}^{(2)}
    -\partial_i g_{j0}^{(3)}
  \right),
  \tag{9.1.18}\label{eq:9.1.18}
\end{equation}
\begin{equation}
  \Gamma^{i(2)}{}_{jk}
  =\frac12\left(
    \partial_k g_{ij}^{(2)}
    +\partial_j g_{ik}^{(2)}
    -\partial_i g_{jk}^{(2)}
  \right),
  \tag{9.1.19}\label{eq:9.1.19}
\end{equation}
\begin{equation}
  \Gamma^{0(3)}{}_{00}=-\frac12\partial_tg_{00}^{(2)},
  \tag{9.1.20}\label{eq:9.1.20}
\end{equation}
\begin{equation}
  \Gamma^{0(2)}{}_{0\ii}=-\frac12\partial_i g_{00}^{(2)},
  \tag{9.1.21}\label{eq:9.1.21}
\end{equation}
\begin{equation}
  \Gamma^{0(1)}{}_{ij}=0.
  \tag{9.1.22}\label{eq:9.1.22}
\end{equation}
Evidently, we shall have to know \(g_{ij}\) to order \(\bar v^2\),
\(g_{i0}\) to order \(\bar v^3\), and \(g_{00}\) to order \(\bar v^4\).
This should be contrasted with the Newtonian approximation, in which we
needed \(g_{00}\) to order \(\bar v^2\), and \(g_{i0}\) and \(g_{ij}\) only
to zeroth order.

To calculate the Ricci tensor we use the target convention
\[
  R_{\mu\nu}
  =
  \partial_\lambda\Gamma^\lambda{}_{\nu\mu}
  -\partial_\nu\Gamma^\lambda{}_{\lambda\mu}
  +\Gamma^\lambda{}_{\lambda\rho}\Gamma^\rho{}_{\nu\mu}
  -\Gamma^\lambda{}_{\nu\rho}\Gamma^\rho{}_{\lambda\mu}.
\]
The components have the expansions
\begin{equation}
  R_{00}=R_{00}^{(2)}+R_{00}^{(4)}+\cdots,
  \tag{9.1.23}\label{eq:9.1.23}
\end{equation}
\begin{equation}
  R_{i0}=R_{i0}^{(3)}+R_{i0}^{(5)}+\cdots,
  \tag{9.1.24}\label{eq:9.1.24}
\end{equation}
\begin{equation}
  R_{ij}=R_{ij}^{(2)}+R_{ij}^{(4)}+\cdots,
  \tag{9.1.25}\label{eq:9.1.25}
\end{equation}
where \(R_{\mu\nu}^{(N)}\) is of order
\(\bar v^N/\bar r^2\). The terms calculable from the ``known'' pieces of the
connection are
\begin{equation}
  R_{00}^{(2)}=\partial_i\Gamma^{i(2)}{}_{00},
  \tag{9.1.26}\label{eq:9.1.26}
\end{equation}
\begin{equation}
  R_{00}^{(4)}
  =
  -\partial_t\Gamma^{i(3)}{}_{0\ii}
  +\partial_i\Gamma^{i(4)}{}_{00}
  -\Gamma^{0(2)}{}_{0\ii}\Gamma^{i(2)}{}_{00}
  +\Gamma^{i(2)}{}_{00}\Gamma^{j(2)}{}_{ij},
  \tag{9.1.27}\label{eq:9.1.27}
\end{equation}
\begin{equation}
  R_{i0}^{(3)}
  =
  -\partial_t\Gamma^{j(2)}{}_{ij}
  +\partial_j\Gamma^{j(3)}{}_{i0},
  \tag{9.1.28}\label{eq:9.1.28}
\end{equation}
\begin{equation}
  R_{ij}^{(2)}
  =
  -\partial_j\Gamma^{0(2)}{}_{i0}
  -\partial_j\Gamma^{k(2)}{}_{ik}
  +\partial_k\Gamma^{k(2)}{}_{ij}.
  \tag{9.1.29}\label{eq:9.1.29}
\end{equation}
Using Eqs.~\eqref{eq:9.1.16}--\eqref{eq:9.1.21}, these become
\begin{equation}
  R_{00}^{(2)}=-\frac12\nabla^2 g_{00}^{(2)},
  \tag{9.1.30}\label{eq:9.1.30}
\end{equation}
\begin{equation}
  \begin{split}
  R_{00}^{(4)}={}&
  -\frac12\partial_t^2g_{ii}^{(2)}
  +\partial_i\partial_tg_{i0}^{(3)}
  -\frac12\nabla^2g_{00}^{(4)}
  +\frac12g_{ij}^{(2)}\partial_i\partial_jg_{00}^{(2)}
  \\
  &+\frac12(\partial_jg_{ij}^{(2)})(\partial_i g_{00}^{(2)})
  -\frac14(\partial_i g_{00}^{(2)})(\partial_i g_{00}^{(2)})
  -\frac14(\partial_i g_{00}^{(2)})(\partial_i g_{jj}^{(2)}),
  \end{split}
  \tag{9.1.31}\label{eq:9.1.31}
\end{equation}
\begin{equation}
  R_{i0}^{(3)}
  =
  -\frac12\partial_i\partial_tg_{jj}^{(2)}
  +\frac12\partial_i\partial_jg_{j0}^{(3)}
  +\frac12\partial_j\partial_tg_{ij}^{(2)}
  -\frac12\nabla^2g_{i0}^{(3)},
  \tag{9.1.32}\label{eq:9.1.32}
\end{equation}
\begin{equation}
  \begin{split}
  R_{ij}^{(2)}={}&
  \frac12\partial_i\partial_jg_{00}^{(2)}
  -\frac12\partial_i\partial_jg_{kk}^{(2)}
  +\frac12\partial_k\partial_jg_{ik}^{(2)}
  \\
  &+\frac12\partial_k\partial_i g_{jk}^{(2)}
  -\frac12\nabla^2g_{ij}^{(2)}.
  \end{split}
  \tag{9.1.33}\label{eq:9.1.33}
\end{equation}

A tremendous simplification can be achieved at this point by choosing a
suitable coordinate system. We showed in Section~\ref{sec:7.4} that it is
always possible to define the \(x^\mu\) so that they obey the harmonic
coordinate conditions
\begin{equation}
  g^{\mu\nu}\Gamma^\lambda{}_{\mu\nu}=0.
  \tag{9.1.34}\label{eq:9.1.34}
\end{equation}
The third-order term with \(\lambda=0\) gives
\begin{equation}
  0=
  \frac12\partial_tg_{00}^{(2)}
  -\partial_i g_{i0}^{(3)}
  +\frac12\partial_tg_{ii}^{(2)},
  \tag{9.1.35}\label{eq:9.1.35}
\end{equation}
while the second-order term with \(\lambda=\ii\) gives
\begin{equation}
  0=
  \frac12\partial_i g_{00}^{(2)}
  +\partial_jg_{ij}^{(2)}
  -\frac12\partial_i g_{jj}^{(2)}.
  \tag{9.1.36}\label{eq:9.1.36}
\end{equation}
It follows, respectively, that
\[
  \frac12\partial_t^2g_{ii}^{(2)}
  -\partial_i\partial_tg_{i0}^{(3)}
  +\frac12\partial_t^2g_{00}^{(2)}=0,
\]
\[
  \partial_i\partial_tg_{jj}^{(2)}
  -\partial_i\partial_jg_{j0}^{(3)}
  -\partial_j\partial_tg_{ij}^{(2)}=0,
\]
\[
  \partial_k\partial_i g_{ij}^{(2)}
  +\partial_i\partial_jg_{kj}^{(2)}
  -\partial_i\partial_k g_{jj}^{(2)}
  +\partial_i\partial_k g_{00}^{(2)}=0.
\]
Equations~\eqref{eq:9.1.30}--\eqref{eq:9.1.33} therefore simplify to
\begin{equation}
  R_{00}^{(2)}=-\frac12\nabla^2g_{00}^{(2)},
  \tag{9.1.37}\label{eq:9.1.37}
\end{equation}
\begin{equation}
  R_{00}^{(4)}
  =
  -\frac12\nabla^2g_{00}^{(4)}
  +\frac12\partial_t^2g_{00}^{(2)}
  +\frac12g_{ij}^{(2)}\partial_i\partial_jg_{00}^{(2)}
  -\frac12\bigl(\nabla g_{00}^{(2)}\bigr)^2,
  \tag{9.1.38}\label{eq:9.1.38}
\end{equation}
\begin{equation}
  R_{0\ii}^{(3)}=-\frac12\nabla^2g_{i0}^{(3)},
  \tag{9.1.39}\label{eq:9.1.39}
\end{equation}
\begin{equation}
  R_{ij}^{(2)}=-\frac12\nabla^2g_{ij}^{(2)}.
  \tag{9.1.40}\label{eq:9.1.40}
\end{equation}

We are now ready to use the Einstein field equations, in trace-reversed form,
\begin{equation}
  R_{\mu\nu}
  =
  8\pi G\left(T_{\mu\nu}-\frac12g_{\mu\nu}T^\lambda{}_\lambda\right).
  \tag{9.1.41}\label{eq:9.1.41}
\end{equation}
From their interpretation as energy density, momentum density, and momentum
flux, we expect
\begin{equation}
  T^{00}=T^{00(0)}+T^{00(2)}+\cdots,
  \tag{9.1.42}\label{eq:9.1.42}
\end{equation}
\begin{equation}
  T^{i0}=T^{i0(1)}+T^{i0(3)}+\cdots,
  \tag{9.1.43}\label{eq:9.1.43}
\end{equation}
\begin{equation}
  T^{ij}=T^{ij(2)}+T^{ij(4)}+\cdots,
  \tag{9.1.44}\label{eq:9.1.44}
\end{equation}
where \(T^{\mu\nu(N)}\) denotes the term in \(T^{\mu\nu}\) of order
\((\bar M/\bar r^3)\bar v^N\). In particular, \(T^{00(0)}\) is the
density of rest mass, while \(T^{00(2)}\) is the nonrelativistic part of
the energy density. What we need is
\begin{equation}
  S_{\mu\nu}\equiv
  T_{\mu\nu}-\frac12g_{\mu\nu}T^\lambda{}_\lambda.
  \tag{9.1.45}\label{eq:9.1.45}
\end{equation}
Since \(G\bar M/\bar r\) is of order \(\bar v^2\),
\begin{equation}
  S_{00}=S_{00}^{(0)}+S_{00}^{(2)}+\cdots,
  \tag{9.1.46}\label{eq:9.1.46}
\end{equation}
\begin{equation}
  S_{i0}=S_{i0}^{(1)}+S_{i0}^{(3)}+\cdots,
  \tag{9.1.47}\label{eq:9.1.47}
\end{equation}
\begin{equation}
  S_{ij}=S_{ij}^{(0)}+S_{ij}^{(2)}+\cdots,
  \tag{9.1.48}\label{eq:9.1.48}
\end{equation}
where \(S_{\mu\nu}^{(N)}\) denotes the term of order
\(\bar M\bar v^N/\bar r^3\). In particular,
\begin{equation}
  S_{00}^{(0)}=\frac12T^{00(0)},
  \tag{9.1.49}\label{eq:9.1.49}
\end{equation}
\begin{equation}
  S_{00}^{(2)}
  =
  \frac12\left[
    T^{00(2)}
    -2g_{00}^{(2)}T^{00(0)}
    +T^{ii(2)}
  \right],
  \tag{9.1.50}\label{eq:9.1.50}
\end{equation}
\begin{equation}
  S_{i0}^{(1)}=-T^{0\ii(1)},
  \tag{9.1.51}\label{eq:9.1.51}
\end{equation}
\begin{equation}
  S_{ij}^{(0)}=\frac12\delta_{ij}T^{00(0)}.
  \tag{9.1.52}\label{eq:9.1.52}
\end{equation}
Using Eqs.~\eqref{eq:9.1.37}--\eqref{eq:9.1.40} and
\eqref{eq:9.1.46}--\eqref{eq:9.1.52} in Eq.~\eqref{eq:9.1.41}, we find
that the field equations in harmonic coordinates are indeed consistent with
our expansions, and give
\begin{equation}
  \nabla^2g_{00}^{(2)}=-8\pi G T^{00(0)},
  \tag{9.1.53}\label{eq:9.1.53}
\end{equation}
\begin{equation}
  \begin{split}
  \nabla^2g_{00}^{(4)}={}&
  \partial_t^2g_{00}^{(2)}
  +g_{ij}^{(2)}\partial_i\partial_jg_{00}^{(2)}
  -(\partial_i g_{00}^{(2)})(\partial_i g_{00}^{(2)})\\
  &-8\pi G\left[
    T^{00(2)}
    -2g_{00}^{(2)}T^{00(0)}
    +T^{ii(2)}
  \right],
  \end{split}
  \tag{9.1.54}\label{eq:9.1.54}
\end{equation}
\begin{equation}
  \nabla^2g_{i0}^{(3)}=16\pi G T^{i0(1)},
  \tag{9.1.55}\label{eq:9.1.55}
\end{equation}
\begin{equation}
  \nabla^2g_{ij}^{(2)}
  =-8\pi G\delta_{ij}T^{00(0)}.
  \tag{9.1.56}\label{eq:9.1.56}
\end{equation}
Equation~\eqref{eq:9.1.53} gives, as expected,
\begin{equation}
  g_{00}^{(2)}=-2\phi,
  \tag{9.1.57}\label{eq:9.1.57}
\end{equation}
where \(\phi\) is the Newtonian potential, defined by Poisson's equation
\begin{equation}
  \nabla^2\phi=4\pi G T^{00(0)}.
  \tag{9.1.58}\label{eq:9.1.58}
\end{equation}
Since \(g_{00}^{(2)}\) must vanish at infinity,
\begin{equation}
  \phi(\mathbf x,t)
  =
  -G\int\frac{T^{00(0)}(\mathbf x',t)}
  {|\mathbf x-\mathbf x'|}\,\dd^3x'.
  \tag{9.1.59}\label{eq:9.1.59}
\end{equation}
Likewise, Eq.~\eqref{eq:9.1.56} gives
\begin{equation}
  g_{ij}^{(2)}=-2\delta_{ij}\phi.
  \tag{9.1.60}\label{eq:9.1.60}
\end{equation}
On the other hand, \(g_{i0}^{(3)}\) is a new vector potential
\(\boldsymbol\zeta\):
\begin{equation}
  g_{i0}^{(3)}=\zeta_i,
  \tag{9.1.61}\label{eq:9.1.61}
\end{equation}
and the solution of Eq.~\eqref{eq:9.1.55} that vanishes at infinity is
\begin{equation}
  \zeta_i(\mathbf x,t)
  =
  -4G\int\frac{T^{i0(1)}(\mathbf x',t)}
  {|\mathbf x-\mathbf x'|}\,\dd^3x'.
  \tag{9.1.62}\label{eq:9.1.62}
\end{equation}

Finally, Eq.~\eqref{eq:9.1.54} may be simplified using
Eqs.~\eqref{eq:9.1.57}, \eqref{eq:9.1.58}, and
\[
  \partial_i\phi\,\partial_i\phi
  =\frac12\nabla^2\phi^2-\phi\nabla^2\phi.
\]
The result is
\begin{equation}
  g_{00}^{(4)}=-2\phi^2-2\psi,
  \tag{9.1.63}\label{eq:9.1.63}
\end{equation}
where \(\psi\) is a second potential satisfying
\begin{equation}
  \nabla^2\psi
  =
  \partial_t^2\phi
  +4\pi G\left[T^{00(2)}+T^{ii(2)}\right].
  \tag{9.1.64}\label{eq:9.1.64}
\end{equation}
Again imposing vanishing at infinity gives
\begin{equation}
  \psi(\mathbf x,t)
  =
  -\int\frac{\dd^3x'}{|\mathbf x-\mathbf x'|}
  \left[
    \frac1{4\pi}\partial_t^2\phi(\mathbf x',t)
    +G T^{00(2)}(\mathbf x',t)
    +G T^{ii(2)}(\mathbf x',t)
  \right].
  \tag{9.1.65}\label{eq:9.1.65}
\end{equation}
The coordinate condition \eqref{eq:9.1.35} imposes the further relation
\begin{equation}
  4\partial_t\phi+\boldsymbol\nabla\cdot\boldsymbol\zeta=0,
  \tag{9.1.66}\label{eq:9.1.66}
\end{equation}
while Eq.~\eqref{eq:9.1.36} is now automatically satisfied. We shall see
in Section~\ref{sec:9.3} that Eq.~\eqref{eq:9.1.66} is also satisfied by
our solutions, by virtue of the conservation conditions obeyed by
\(T^{\mu\nu}\).

Using Eqs.~\eqref{eq:9.1.57}, \eqref{eq:9.1.60},
\eqref{eq:9.1.61}, and \eqref{eq:9.1.63} in
Eqs.~\eqref{eq:9.1.16}--\eqref{eq:9.1.22} gives the desired components
of the affine connection:
\begin{equation}
  \Gamma^{i(2)}{}_{00}=\partial_i\phi,
  \tag{9.1.67}\label{eq:9.1.67}
\end{equation}
\begin{equation}
  \Gamma^{i(4)}{}_{00}
  =
  \partial_i(2\phi^2+\psi)+\partial_t\zeta_i,
  \tag{9.1.68}\label{eq:9.1.68}
\end{equation}
\begin{equation}
  \Gamma^{i(3)}{}_{0j}
  =
  \frac12(\partial_j\zeta_i-\partial_i\zeta_j)
  -\delta_{ij}\partial_t\phi,
  \tag{9.1.69}\label{eq:9.1.69}
\end{equation}
\begin{equation}
  \Gamma^{i(2)}{}_{jk}
  =
  -\delta_{ij}\partial_k\phi
  -\delta_{ik}\partial_j\phi
  +\delta_{jk}\partial_i\phi,
  \tag{9.1.70}\label{eq:9.1.70}
\end{equation}
\begin{equation}
  \Gamma^{0(3)}{}_{00}=\partial_t\phi,
  \tag{9.1.71}\label{eq:9.1.71}
\end{equation}
\begin{equation}
  \Gamma^{0(2)}{}_{0\ii}=\partial_i\phi.
  \tag{9.1.72}\label{eq:9.1.72}
\end{equation}
As a bonus, we can calculate three additional terms that will play a role
in post-Newtonian hydrodynamics:
\begin{equation}
  \Gamma^{0(3)}{}_{ij}
  =
  -\frac12(\partial_i\zeta_j+\partial_j\zeta_i)
  -\delta_{ij}\partial_t\phi,
  \tag{9.1.73}\label{eq:9.1.73}
\end{equation}
\begin{equation}
  \Gamma^{0(4)}{}_{i0}=\partial_i\psi,
  \tag{9.1.74}\label{eq:9.1.74}
\end{equation}
\begin{equation}
  \Gamma^{0(5)}{}_{00}
  =
  \partial_t\psi+\boldsymbol\zeta\cdot\boldsymbol\nabla\phi.
  \tag{9.1.75}\label{eq:9.1.75}
\end{equation}
